\documentclass[11pt]{article}
\usepackage{graphicx} % Required for inserting images
\usepackage{amsmath}
\usepackage{amssymb}
\usepackage{xcolor}
\usepackage{float}
\usepackage{listings}
\usepackage{xcolor}

\lstset{
    % change caption to be below the listing
    % choose appropriate font style.
    captionpos=b,
  basicstyle=\footnotesize\ttfamily,
  numbers=left,
  numberstyle=\tiny,
  stepnumber=1,
  numbersep=8pt,
  frame=single,
  breaklines=true,
  tabsize=2,
  showstringspaces=false,
  keywordstyle=\color{blue},
  commentstyle=\color{gray},
  stringstyle=\color{teal},
}

\title{Shussman \& Heizler (2015) - Summary \& Questions}
\author{Tal Ramot}
\date{November 2025}

\begin{document}

\maketitle

\section{Introduction}
\subsection{Statement of the problem}

\section{II Supersonic Waves}
Eq. (9) in this article is the equation over the (only) dimensionless profile of $T(\xi)$ received by substituting the relations for $\frac{d\tilde{T}}{dt}$ and $\frac{d\tilde{T}}{dm}$ using the definition of the similarity dimensionless variable:
\begin{equation}
  \xi = mA^{\frac{1}{2}} T_0^{\frac{\beta - \alpha - 4}{2}} t^{\frac{\tau(\beta - \alpha - 4)-1}{2}}
\end{equation} 
The equation is:
\begin{equation}
  \tilde{T}^{\beta - \alpha - 4} \left[\tau \tilde{T} + w_3 \xi \frac{d\tilde{T}}{d\xi}\right] = \left(\alpha+3\right)\tilde{T}^{-1} \left(\frac{d\tilde{T}}{d\xi}\right)^2 + \frac{\partial^2\tilde{T}}{\partial \xi^2}
\end{equation}
solving this eq. (9) for the second derivative of $\tilde{T}$ gives:
\begin{equation}
  \frac{\partial^2\tilde{T}}{\partial \xi^2} = \tilde{T}^{\beta - \alpha - 4} \left[\tau \tilde{T} + w_3 \xi \frac{d\tilde{T}}{d\xi}\right] - \left(\alpha+3\right)\tilde{T}^{-1} \left(\frac{d\tilde{T}}{d\xi}\right)^2
\end{equation}
One can convert the 2nd order ODE to a system of two 1st order ODEs by defining $y_1 = \tilde{T}$ and $y_2 = \frac{d\tilde{T}}{d\xi}$, which gives the system of coupled 1st order ODEs:
\begin{equation}
  \frac{dy_1}{d\xi} = y_2
\end{equation}
\begin{equation}
  \frac{dy_2}{d\xi} = y_1^{\beta - \alpha - 4} \left[\tau y_1 + w_3 \xi y_2\right] - \left(\alpha+3\right)y_1^{-1} y_2^2
\end{equation}
which is exactly what Shussman solver does:
\begin{lstlisting}[language=Matlab]
  function xp=F(t,x,alpha,beta,tau)
  W3=-0.5+0.5*(beta-alpha-4)*tau;
  xp=zeros(2,1);
  xp(1)=x(2);
  xp(2)=(x(1)^(beta-alpha-4))*[W3*t*x(2)+tau*x(1)]-(alpha+3)*(1/x(1))*(x(2)^2);
  end
\end{lstlisting}
To sum up, the equation that is being solved here is:
\begin{equation}
  \frac{d}{d\xi}
  \begin{bmatrix}
    y_1 \\
    y_2
    \end{bmatrix} =
    \begin{bmatrix}
    y_2 \\
    y_1^{\beta - \alpha - 4} \left[\tau y_1 + w_3 \xi y_2\right] - \left(\alpha+3\right)y_1^{-1} y_2^2
    \end{bmatrix}
\end{equation}
which can be solved numerically using the Runge-Kutta method which is implemented by the \texttt{ode45} function in Matlab.
\subsection{The initial and boundary conditions \& Shooting method}
In order to solve this system of coupled ODEs, one should specify a well-defined initial condition. However, the one wishes to find the dimensionless temperature profile $y_1 = \tilde{T}$ that satisfies the following \textbf{Boundary Conditions}:
\begin{equation}
  \begin{gathered}
    \tilde{T}(\xi \to \xi_f) = 0 \\
  \frac{d\tilde{T}}{d\xi}(\xi \to \xi_f) = -\infty
  \end{gathered}
\end{equation}
with the normalization condition:
\begin{equation}
  \tilde{T}(\xi = 0) = 1
\end{equation}
Therefore, the solution is solved numerically using the \textbf{Shooting method}, which is a numerical method for solving boundary value problems by converting them into initial value problems (IVP). Therefore, one can solve this system of coupled ODEs by guessing the initial value of:
\begin{equation}
  \begin{aligned}
    \begin{bmatrix}
      y_1(\xi_f) \\
      y_2(\xi_f)
    \end{bmatrix} = 
    \begin{bmatrix}
      0 \\
      -1000
    \end{bmatrix}
  \end{aligned}
\end{equation}
for an arbitrary value of $\xi_f$ (the front dimensionless position), and then integrate the equation from $xi_f$ to $0$.
\section{The patching method}
In case of a general problem of a power-law temperature drive at the boundary, a solution of the hydrodynamic \& radiation profiles at the entire domain can be solved through dividing the domain into two parts: 


\section{Questions}
\textbf{Questions to answer - Week 13 (15.1.26)}
\begin{enumerate}
    \item Understand the introduction - what exactly is the innovation shown in this paper? 
    \item Understand the equation in subsection II.A.
    \item Understand the self-similarity of the first kind - subsection II.B.
    \item Why the units of the temperature can be defined by
    forcing $f\rho_0^{-\mu}$ to be dimensionless. The problem can also be solved using four units and f as an additional dimensional
    constant.
    \item Understand the physical differrence between the supersonic and subsonic cases - what is the affect of the sonic point?
    \\ \dots
    \item Understand why is there an additional dimensional constant which is $\chi_f$ (for constant and steep heat profile at the front) and/or $\bar{\chi}$ for a known observed quantity. 
    \item How one derived Eq. 18?
    \item Show myself that I can get Eq. 20.
    \item \textcolor{red}{Understand the physical process of the conformation of the Marshak wave!}
    \item \textcolor{red}{Understand the intuition about self-similarity, what is the time-scale in these questions? I}
    \item In general, I try to understand the big picture - it feels like the unified Krief article summarizes the theory that Krief presents... What is better - study his articles or the NLD lecture notes? I think I should do both - but it will take time, I want to choose the best option.
\end{enumerate}

\end{document}


